%% Esakia spaces, Heyting homomorphisms, and the order back condition.

\subsection{Esakia duality}
\label{sec:p10-esakia-duality}

\begin{hbdefinition}{def:p10-esakia-space}{Esakia spaces}
For a subset \(S\) of a partially ordered set, put
\[
 \downarrow S=\{x:\exists s\in S\ (x\leq s)\},\qquad
 \uparrow x=\{y:x\leq y\}.
\]
An Esakia space is a Priestley space \(X\) such that
\(\downarrow C\) is clopen for every clopen subset \(C\) of \(X\).
The condition applies to all clopens, including those that are
not upsets.
The empty Esakia space is allowed, corresponding to the trivial
Heyting algebra.
For a Heyting algebra \(H\), the underlying order of \(X_H\) is
\((\operatorname{Pf}(H),\subseteq)\) from
Definition~\ref{def:p10-heyting-prime-frame}.
As point sets, \(\eta_H(a)=\beta_H(a)=\widehat a\); the Priestley
topology identifies these sets as clopen upsets.
This is the modern Priestley-space formulation in
\cite[Definition~2.3.20, p.~34]{Bezhanishvili2006Thesis}.
\index{Esakia space}\index{Esakia duality}
\end{hbdefinition}

\begin{hbtheorem}{thm:p10-esakia-objects}{proved=annex:p10.esakia.1}{Heyting algebras and Esakia spaces}
Assume BPIT.
The Priestley space \(X_H\) of a Heyting algebra \(H\) is Esakia,
and
\[
 \downarrow(\widehat a\setminus\widehat b)
       =X_H\setminus\widehat{a\Rightarrow b}.
\]
Conversely, the clopen upsets of an Esakia space form a Heyting
algebra with
\[
 U\Rightarrow V=X\setminus\downarrow(U\setminus V).
\]
The map \(\eta_H:a\mapsto\widehat a\) is a Heyting isomorphism
onto \(\operatorname{ClopUp}(X_H)\).
The evaluation map
\(x\mapsto\{U\in\operatorname{ClopUp}(X):x\in U\}\)
recovers \(X\) as an ordered topological space.
\end{hbtheorem}

\begin{hbdefinition}{def:p10-esakia-morphism}{Esakia morphisms}
An Esakia morphism \(g:X\to Y\) is a continuous order-preserving
map satisfying the order back condition
\[
 g(x)\leq y\quad\Longrightarrow\quad
       \exists z\geq x\quad g(z)=y.
\]
Equivalently, \(g[\uparrow x]=\uparrow g(x)\) for every \(x\).
Order preservation gives one inclusion and back gives the other.
These are continuous p-morphisms of the underlying partial orders
\cite[\S2.3.3, pp.~35--37]{Bezhanishvili2006Thesis}.
\end{hbdefinition}

\begin{hblemma}{thm:p10-heyting-prime-lift}{proved=annex:p10.esakia.2}{Lifting prime filters along a Heyting homomorphism}
Assume BPIT.
Let \(h:H\to K\) be a Heyting homomorphism, \(F\) a prime filter
of \(K\), and \(G\) a prime filter of \(H\), with
\(h^{-1}[F]\subseteq G\).
There is a prime filter \(F'\supseteq F\) of \(K\) satisfying
\[
 h^{-1}[F']=G.
\]
Thus the inverse-image map \(X_K\to X_H\) has the order back
condition.
\end{hblemma}

\begin{hbtheorem}{thm:p10-esakia-functors}{proved=annex:p10.esakia.3}{Functorial Esakia duality}
Assume BPIT.
Heyting algebras with Heyting homomorphisms are contravariantly
equivalent to Esakia spaces with Esakia morphisms.
The functors are the prime-filter and clopen-upset constructions.
On arrows they act by inverse image:
\[
\begin{aligned}
 h:H\to K&\quad\longmapsto\quad
     X(h):X_K\to X_H,\quad F\mapsto h^{-1}[F],\\
 g:X\to Y&\quad\longmapsto\quad
     \operatorname{ClopUp}(g):
     \operatorname{ClopUp}(Y)\to\operatorname{ClopUp}(X).
\end{aligned}
\]
The Priestley reconstruction isomorphisms are natural for these
restricted objects and arrows.
The trivial algebra and empty space are included.
\end{hbtheorem}

Esakia's original paper formulates the duality using perfect
topological Kripke models and strongly isotone maps
\cite[Definition~2, p.~298; Theorem~3, case \(\alpha=1\), p.~300]{Esakia1974}.
Its representation on p.~301 uses clopen upper cones, with
implication determined by their future-point condition.
The explicit modern prime-filter and Priestley-space
presentation above follows
\cite[Theorems~2.3.22 and~2.3.25, pp.~35, 37]{Bezhanishvili2006Thesis};
the annex supplies the full filter-lifting and morphism proofs.
The original paper and that thesis use nonempty spaces.
The empty-space case is included explicitly here to retain
the trivial Heyting algebra.

\paragraph{A monotone map need not preserve implication dually.}
Let \(X=\{x\}\) and \(Y=\{0<1\}\), both with the discrete topology.
These finite ordered spaces are Esakia.
The map \(g(x)=0\) is continuous and order-preserving, but does
not satisfy back, since \(1\geq g(x)\) has no preimage.
For the clopen upsets \(U=\{1\}\), \(V=\varnothing\) of \(Y\),
\[
 U\Rightarrow V=\varnothing,\qquad
 g^{-1}[U]\Rightarrow g^{-1}[V]=X.
\]
Thus inverse image under this map does not preserve Heyting
implication. The order back condition is part of the arrow
correspondence.
